% =====================================================================
%  config/dados.tex  --  DADOS DO SEU TRABALHO (EDITE ESTE ARQUIVO!)
% ---------------------------------------------------------------------
%  Tudo o que você escrever aqui aparece automaticamente na capa, na
%  folha de rosto e na folha de aprovação. Basta trocar o texto que está
%  entre as chaves { }. Não apague as chaves nem a barra invertida (\).
% =====================================================================

% --- Instituição e curso ---------------------------------------------
\instituicao{Centro Universitário Ibmec}

% Escolha UM dos cursos abaixo (deixe os outros com % na frente):
\newcommand{\imprimircurso}{Engenharia de Computação}
% \newcommand{\imprimircurso}{Engenharia Civil}
% \newcommand{\imprimircurso}{Engenharia Mecânica}
% \newcommand{\imprimircurso}{Engenharia de Produção}

% Grau pretendido (normalmente não precisa mudar)
\newcommand{\imprimirgrau}{Bacharel em \imprimircurso}

% --- Autor e trabalho -------------------------------------------------
% Escreva seu nome completo normalmente (maiúsculas e minúsculas).
% A capa converte para CAIXA ALTA sozinha.
\autor{André Costa Coelho}

% Título: claro e preciso. Se houver subtítulo, use o comando \subtitulo.
\titulo{RoundUp}
% Subtítulo (opcional). Deixe vazio {} se não houver.
\newcommand{\imprimirsubtitulo}{aplicativo gamificado para engajamento e retenção de praticantes de artes marciais}

% Título em inglês (usado no Abstract)
\tituloestrangeiro{RoundUp: a gamified app for engagement and retention of martial arts practitioners}

% --- Local e ano ------------------------------------------------------
\local{Rio de Janeiro}
\data{2026}

% --- Orientação -------------------------------------------------------
% Titulações (Manual, 2.1.4): Especialista = Esp.; Mestre = Ms.; Doutor = Dr.
% Para orientadora, use "Profa." e "Orientadora".
\orientador{Profa. Dra. Talita Vieira Ribeiro}  % confira a titulação (Ms./Dra.)
\newcommand{\imprimirrotuloorientador}{Orientadora}  % ou Orientadora
% Coorientador (opcional): retire o % da linha abaixo
% \coorientador{Prof. Ms. Nome do Coorientador}

% --- Banca examinadora (Regulamento, Art. 14) ------------------------
% Membro 1 = orientador (preside a banca); membros 2 e 3 = examinadores.
\newcommand{\bancaorientador}{Profa. Dra. Talita Vieira Ribeiro}
\newcommand{\bancaorientadorinst}{Centro Universitário Ibmec}
\newcommand{\bancaexaminadorum}{Prof. Ms. Nome do Examinador 1}
\newcommand{\bancaexaminadoruminst}{Centro Universitário Ibmec}
\newcommand{\bancaexaminadordois}{Prof. Dr. Nome do Examinador 2}
\newcommand{\bancaexaminadordoisinst}{Instituição do Examinador 2}

% Data da defesa (deixe em branco até a defesa)
\newcommand{\dataaprovacao}{\underline{\hspace{1cm}}/\underline{\hspace{1cm}}/\underline{\hspace{1.5cm}}}

% --- Informações do PDF (aparecem nas propriedades do arquivo) --------
\makeatletter
\AtBeginDocument{\hypersetup{
  pdftitle={\@title},
  pdfauthor={\@author},
  pdfsubject={Trabalho de Conclusão de Curso -- \imprimircurso},
  pdfcreator={LaTeX com abnTeX2},
}}
\makeatother
